\documentclass[10pt,a4paper]{amsart}
\usepackage[margin=19mm]{geometry}
\usepackage{amsmath,amssymb,mathtools}
\usepackage[scr=rsfs,cal=euler]{mathalfa}
\usepackage{xfrac}
\usepackage[unicode,hidelinks,bookmarks=false]{hyperref}
\newcommand{\ie}{i.e.\ }
\newcommand{\cf}{cf.\ }
\newcommand{\resp}{respectively\ }
\newcommand{\Subsection}{\subsection}
\providecommand{\nobreakdash}{\nobreak\hbox{-}}
\setlength{\emergencystretch}{2em}
\multlinegap0pt
\usepackage{mathrsfs}
\usepackage{tikz-cd}
\usepackage{todonotes}
\newtheorem{thm}{Theorem}[section]
\newtheorem{prop}[thm]{Proposition}
\newtheorem{lem}[thm]{Lemma}
\newtheorem{cor}[thm]{Corollary}
\newtheorem{mainthm}{Theorem}
\newtheorem{defn}[thm]{Definition}
\newtheorem{rem}[thm]{Remark}
\newtheorem{ex}[thm]{Example}
\renewcommand{\themainthm}{\Alph{mainthm}}
\newcommand{\abs}[1]{\left\vert#1\right\vert}
\newcommand{\OO}{\mathcal{O}}
\newcommand{\C}{\mathbb{C}}
\newcommand{\PP}{\mathbb{P}}
\newcommand{\Z}{\mathbb{Z}}
\newcommand{\R}{\mathbb{R}}
\DeclareMathOperator{\sHom}{\mathscr{H}\!om}
\DeclareMathOperator{\sExt}{\mathscr{E}\!xt}
\DeclareMathOperator{\sEnd}{\mathscr{E}\!nd}
\DeclareMathOperator{\Ext}{Ext}
\DeclareMathOperator{\coker}{coker}
\DeclareMathOperator{\im}{im}
\DeclareMathOperator{\codim}{codim}
\DeclareMathOperator{\Spec}{Spec}
\DeclareMathOperator{\Supp}{Supp}
\DeclareMathOperator{\Pic}{Pic}
\DeclareMathOperator{\id}{id}
\DeclareMathOperator{\Bl}{Bl}
\newcommand{\ev}{\operatorname{ev}}      % evaluation map
\newcommand{\Mbar}{\overline{\mathcal M}} % stable maps moduli
\newcommand{\vT}{\mathbb{T}}           % Global deformation groups T^i
\newcommand{\I}{\mathscr{I}}          % Ideal sheaf
\newcommand{\derR}{\mathbf{R}}        % Derived functor R
\newcommand{\push}{\pi_*}             % Standard pushforward
\newcommand{\drpush}{\derR\pi_*}      % Total derived pushforward
\newcommand{\can}{\omega}             % Canonical bundle
\newcommand{\dualV}{^{\vee}}          % Dual (transpose/dual)
\newcommand{\tors}{\text{tors}}       % Torsion
\begin{document}
\title[Log Conifold Transitions]{Log Conifold Transitions}
\author{Rodolfo Aguilar}
\date{}
\maketitle
\noindent\emph{Source and licence.} Log Conifold Transitions. Version arXiv:2606.31072v1 (CC BY 4.0). This material is distributed under the Creative Commons Attribution 4.0 International licence, http://creativecommons.org/licenses/by/4.0/, as recorded in the arXiv OAI licence metadata for this exact version and shown on the abstract page https://arxiv.org/abs/2606.31072v1. Full rights notice: licenses/arxiv-2606.31072-CC-BY-4.0.txt.\par\smallskip
\noindent\textit{Editorial scope.} Counted unit: the original \texttt{thm} environment \texttt{thm:unobstructedness} at source lines 350--358 together with its complete proof at lines 359--389; the delivered document keeps source lines 171--390 so that every referenced label and every citation resolves inside it. Native editable source is the paper's own concatenated TeX; the AMS wrapper is editorial formatting only. No publisher class, upstream script or unrelated artwork is redistributed.
\par\smallskip\noindent\textit{Reference note.} This article ships no compiled bibliography (biblatex); the entries delivered here are composed from the author's own BibTeX (.bib) fields, with no field invented and no entry dropped.
\subsection{Deformation of threefolds and pairs}\label{ss:DeformationPairs}
Here we collect some results adapting the deformation theory of \cite{F86} to the logarithmic setting. 

Denote by:
\begin{enumerate}
  \item $X_0$ a compact analytic threefold.
  \item $Z=\{p_1,\ldots, p_j\}$ the singular locus of $X_0$, which consists of isolated Gorenstein singularities admitting a small resolution.
  \item $\pi: \tilde{X}\to X_0$ a small resolution of $X_0$ at all the points in $Z$.
  \item $\Gamma=\pi^{-1}(Z)$ the exceptional locus, which consists of a union of rational curves.
  \item $Y_0\subset X_0$ a smooth divisor such that $Y_0 \cap Z = \emptyset$.
\end{enumerate}

Because the divisor $Y_0$ avoids the singular locus $Z$ (i.e., $Y_0 \cap Z = \emptyset$), we can define the global sheaf of logarithmic K\"ahler differentials $\Omega_{X_0}^1(\log Y_0)$ without requiring generalized cotangent complexes for singular spaces. We do this by defining the sheaf on an open cover, or equivalently, directly on its stalks.

Consider the Zariski open cover of $X_0$ given by $U = X_0 \setminus Z$ (the smooth locus) and $V = X_0 \setminus Y_0$. 
On $U$, the space is smooth and $Y_0 \subset U$ is a smooth divisor, so the classical Deligne sheaf of logarithmic differentials $\Omega_U^1(\log Y_0)$ is well-defined. 
On $V$, there is no boundary divisor, so we utilize the standard sheaf of K\"ahler differentials $\Omega_V^1$. 
On the intersection $U \cap V = X_0 \setminus (Z \cup Y_0)$, the space is smooth and the boundary is empty. Consequently, both local sheaves canonically restrict to the regular cotangent sheaf $\Omega_{U \cap V}^1$. Because they naturally agree on the overlap, they glue to define a unique global coherent sheaf $\Omega_{X_0}^1(\log Y_0)$ on $X_0$.

Equivalently, this sheaf is completely characterized by its stalks at any point $x \in X_0$:
$$
\left( \Omega_{X_0}^1(\log Y_0) \right)_x = 
\begin{cases} 
\Omega_{X_0, x}^1(\log Y_0) & \text{if } x \in Y_0 \text{ (where } X_0 \text{ is guaranteed to be smooth)}, \\
\Omega_{X_0, x}^1 & \text{if } x \notin Y_0 \text{ (which includes all singular points } p_i \in Z).
\end{cases}
$$
By construction, because $\Omega_{X_0}^1(\log Y_0)$ is isomorphic to the standard K\"ahler differentials in an open neighborhood of $Z$, the local deformation sheaves $T_{X_0}^i(-\log Y_0) := \mathcal{E}xt^i(\Omega_{X_0}^1(\log Y_0), \OO_{X_0})$ are unaffected by the boundary $Y_0$ near the singular locus.
Following Friedman's notation, let $T_{X_0}^0(-\log Y_0) = \mathcal{H}om(\Omega_{X_0}^1(\log Y_0), \OO_{X_0})$ be the logarithmic tangent sheaf, and define the local deformation sheaves of the pair as:
\[
T_{X_0}^i(-\log Y_0) := \mathcal{E}xt^i(\Omega_{X_0}^1(\log Y_0), \OO_{X_0}).
\]
Because $\Omega_{X_0}^1(\log Y_0)$ is locally free away from $Z$ (since $Y_0$ is smooth) and isomorphic to $\Omega_{X_0}^1$ away from $Y$, the higher local Ext sheaves are supported strictly on $Z$. Consequently, for $i > 0$, the boundary condition vanishes locally, yielding a canonical isomorphism:
\[
T_{X_0}^i(-\log Y_0) \cong T_{X_0}^i \cong \bigoplus_{k=1}^j T_{p_k}^i
\]
where $T_{p_k}^i$ is the local deformation space of the singularity $p_k$.

The global deformation spaces of the pair, denoted $\mathbb{T}_{X_0}^i(-\log Y_0) = \operatorname{Ext}^i(\Omega_{X_0}^1(\log Y_0), \OO_{X_0})$, are related to the local deformations via the local-to-global spectral sequence:
\[
E_2^{p,q} = H^p(X_0, T_{X_0}^q(-\log Y_0)) \Rightarrow \mathbb{T}_{X_0}^{p+q}(-\log Y_0).
\]
Similarly, the deformations of the log resolution $(\tilde{X}, \tilde{Y})$ are governed by $H^i(\tilde{X}, T_{\tilde{X}}(-\log \tilde{Y}))$. The Leray spectral sequence for the resolution $\pi$ yields:
\[
E_2^{p,q} = H^p(X_0, R^q \pi_* T_{\tilde{X}}(-\log \tilde{Y})) \Rightarrow H^{p+q}(\tilde{X}, T_{\tilde{X}}(-\log \tilde{Y})).
\]
Because $\pi$ is an isomorphism outside $Z$ and $\tilde{Y}$ avoids $\Gamma$, there is a natural identification $\pi_* T_{\tilde{X}}(-\log \tilde{Y}) \cong T_{X_0}(-\log Y)$, allowing us to compare the $E_2$ pages of these spectral sequences directly.

\subsection{Local Cohomology of the Log Tangent Sheaf}
Let $(X_0,0)$ be the germ of an isolated Gorenstein threefold singularity, and $\pi:\tilde{X}\to X_0$ a small resolution. Denote the exceptional locus by 
$$\pi^{-1}(0)=\cup_i C_i=\Gamma,$$
where the $C_i$ are the irreducible curves. 
Recall that by definition of $\Omega_{X_0}^1(\log Y_0)$, it coincides with $\Omega_{X_0}^1$ around a point $0\in X_0$ away from $Y_0$.
\begin{lem}[{\cite[Lemma 3.1]{F86}}]\label{lem:log_tangent_pushforward} $R^0\push T_{\tilde{X}}(-\log \tilde{Y})\cong T^0_{X_0}(-\log Y_0)\cong T_{X_0}^0 .$
\end{lem}

\begin{lem}[{\cite[Lemma 3.2]{F86}}] \label{lem:log_higher_direct_images}
The local cohomology group $H^2_{\{0\}}(R^1\push T_{\tilde{X}}(-\log \tilde{Y}))=0$
\end{lem}

\begin{lem}[{\cite[Lemma 3.3]{F86}}]\label{lem:Schlessinger} $H^0(T_{X_0}^1(-\log Y_0))=H^2_{\{0\}}(T_{X_0}^0(-\log Y_0))=H^2_{\{0\}}(T_{X_0}^0) $  
\end{lem}

\subsection{The Local-to-Global Spectral Sequences}
Let us now come back to the global setting and use the notation of subsection \ref{ss:DeformationPairs}.

\begin{prop}
  The Leray spectral sequence for the log resolution $H^p(R^q \push T_{\tilde{X}}(-\log \tilde{Y}))\Rightarrow H^{p+q}(T_{\tilde{X}}(-\log \tilde{Y}))$ and the local-to-global spectral sequence for the singular pair $\mathbb{T}_{X_0}^i(-\log Y_0)$ induce the following commutative diagram with exact rows:
\begin{equation} \label{eq:friedman_diagram}
\resizebox{\displaywidth}{!}{%
\begin{tikzcd}[column sep=1.5em, row sep=3em, ampersand replacement=\&]
0 \arrow[r] \& H^1( T^0_{X_0}(-\log Y_0)) \arrow[r] \arrow[d, "\cong"] \& H^1(T_{\tilde{X}}(-\log \tilde{Y})) \arrow[r] \arrow[d, "\phi"] \& H^0(R^1\push T_{\tilde{X}}(-\log \tilde{Y})) \arrow[r, "ob"] \arrow[d] \& H^2(T^0_{\tilde{X}}(-\log \tilde{Y})) \arrow[r] \arrow[d, "\cong"] \& H^2(T_{\tilde{X}}(-\log \tilde{Y})) \arrow[d] \arrow[r] \& 0 \\
0 \arrow[r] \& H^1( T^0_{X_0}(-\log Y_0)) \arrow[r] \& \vT^1_{X_0}(-\log Y_0) \arrow[r] \& H^0(X_0, T^1_{X_0}(-\log Y_0)) \arrow[r, "ob"] \& H^2(T^0_{X_0}(-\log Y_0)) \arrow[r] \& \vT^2_{X_0}(-\log Y_0) \arrow[r] \& 0
\end{tikzcd}%
}
\end{equation}
\end{prop}
\begin{proof}
The top row is the standard five-term exact sequence of low degrees associated to the Leray spectral sequence $E_2^{p,q} = H^p(X_0, R^q\push T_{\tilde{X}}(-\log \tilde{Y})) \Rightarrow H^{p+q}(\tilde{X}, T_{\tilde{X}}(-\log \tilde{Y}))$. Because $R^2\push = 0$, the $E_2^{0,2}$ term vanishes, allowing the sequence to terminate exactly at $H^2(\tilde{X}, T_{\tilde{X}}(-\log \tilde{Y}))$.

The bottom row is the analogous five-term exact sequence for the local-to-global spectral sequence $E_2^{p,q} = H^p(X_0, T^q_{X_0}(-\log Y_0)) \Rightarrow \vT^{p+q}_{X_0}(-\log Y_0)$. Because ordinary double points are local complete intersections, $T^2_{X_0} = 0$, allowing the sequence to terminate at $\vT^2_{X_0}(-\log Y_0)$. 

Furthermore, because the singularities are isolated points and the resolution is small, both $R^1\push T_{\tilde{X}}(-\log \tilde{Y})$ and $T^1_{X_0}(-\log Y_0)$ are skyscraper sheaves supported entirely on the zero-dimensional locus $Z$. Consequently, their first cohomology groups vanish 
($H^1(X_0, R^1\push T_{\tilde{X}}(-\log \tilde{Y}))=0$ and $H^1(X_0, T^1_{X_0}(-\log Y_0))=0$), which forces the maps originating from the degree-two base spaces to be surjective.
\end{proof}


\section{Unobstructedness}\label{s:Un}
\subsection{Setting and basic sequences}\label{ss:setting}
Let us fix the setting and notation for this section. Let $X$ be a smooth prime Fano threefold of index $2$ and degree $d\in \{2,3,4,5,8\}$.
Let $Y\in \abs{-\frac{1}{2} K_X}$ be a generic smooth divisor and $\xi\subset Y$ a generic subset of $e$ points.
Denote by $\sigma:\tilde{X}\to X$ the blow-up of $X$ at $\xi$ and by $\tilde{Y}$ the strict transform of $Y$, by $\Gamma=\cup \tilde{C}_i\subset \tilde X\setminus \tilde Y$ a union of smooth rational curves assumed \emph{disjoint} with normal bundle $N_{\tilde C_i/\tilde X}\cong \OO(-1)\oplus \OO(-1)$.
Let $\pi:\tilde{X}\to X_0$ be the contraction of $\Gamma$ with $\pi_*(\Gamma)=Z$ and $Y_0\cong \pi_*(\tilde Y)$.

For deformations of the threefold $X_0$ we follow the notation of \cite{F86}. 
As for deformations of the pair $(X_0,Y_0)$ we use the notation and definitions of section \ref{ss:DeformationPairs}. 

There is the exact sequence
$$0\to T_{X_0}^0(-\log Y_0)\to T_{X_0}^0\to N_{Y_0/X_0}\to 0. $$
As $N_{Y_0/X_0}=K_{Y_0}^{-1}$, $H^2(Y_0, N_{Y_0/X_0})=H^2(Y_0, K_{Y_0}^{-1})$ is Serre dual to $H^0(Y_0, K_{Y_0}^2)=0$. 
So there is an exact sequence
\begin{equation}\label{eq:LEStanToNor}
  H^1(X_0, T_{X_0}^0(-\log Y_0))\to H^1(X_0, T_{X_0}^0)\to H^1(Y_0, N_{Y_0/X_0})\to H^2(X_0, T_{X_0}^0(-\log Y_0))\to 0. 
\end{equation}
The local-to-global $\Ext$ spectral sequence gives:
{\small 
\begin{equation} \label{eq:locToGlo} % Changed label name
\begin{gathered}
\begin{tikzcd}[column sep=0.5em] % Shrunk from 1em to 0.5em to fit page
\Ext^1(\Omega_{X_0}^1(\log Y_0), \OO_{X_0}) \ar[d, "\cong "] \ar[r] & H^0(\sExt^1(\Omega_{X_0}^1(\log Y_0), \OO_{X_0}))\ar[d, "\cong "] \ar[r]  & H^2(X_0, T_{X_0}^0(-\log Y_0)) \ar[r] & \Ext^2(\Omega_{X_0}^1(\log Y_0), \OO_{X_0}) \ar[r] \ar[d, "\cong"] & 0 \\
\mathbb{T}_{X_0}^1(-\log Y_0)  &  H^0(T_{X_0}^1) &      & \mathbb{T}_{X_0}^2(-\log Y_0) & 
\end{tikzcd}
\end{gathered}
\end{equation}
}
Similarly, the Poincaré residue exact sequence 
$$0\to \Omega_{X_0}^1\to \Omega_{X_0}^1(\log Y_0)\to \OO_{Y_0}\to 0, $$
gives the long exact $\Ext$ sequence
\begin{equation}\label{eq:DefPairToDef}
  \begin{tikzcd}
  \Ext^1(\Omega_{X_0}^1(\log Y_0),\OO_{X_0}) \ar[d, "\cong "] \ar[r] & \Ext^1(\Omega_{X_0}^1,\OO_{X_0}) \ar[d, "\cong"] \ar[r] & \Ext^2(\OO_{Y_0}, \OO_{X_0})\ar[d, "\cong"] \ar[r] & 0\\
  \mathbb{T}_{X_0}^1(-\log Y_0) & \mathbb{T}_{X_0}^1 & H^1(Y_0,N_{Y_0/X_0}) &
\end{tikzcd}.
\end{equation}
Here $\Ext^2(\OO_{Y_0}, \OO_{X_0})\cong H^1(\sExt^1(\OO_{Y_0}, \OO_{X_0}))\cong H^1(\OO_{X_0}(Y_0)/\OO_{X_0})\cong H^1(Y_0, N_{Y_0/X_0}).$
\subsection{Deformations of the threefold}

\begin{prop}\label{prop:defAbs} The deformations of $X_0$ and of $\tilde X$ are unobstructed, this is, 
  \begin{enumerate}
  \item $\mathbb{T}_{X_0}^2=0$,
  \item $H^2(X_0,T_{X_0}^0)\cong H^2(\tilde X, T_{\tilde X})=0$.
  \end{enumerate} 
  As a consequence, the map $\mathbb{T}_{X_0}^1\to H^0(X_0, T_{X_0}^1)$ is surjective.
\end{prop}
\begin{proof}
  We have the following exact sequence given by the local-to-global spectral sequence:
  $$\mathbb{T}_{X_0}^1\to H^0(X_0,T_{X_0}^1)\to H^2(X_0, T_{X_0}^0)\to \mathbb{T}_{X_0}^2\to 0. $$
  Because $Z$ consists of ordinary double points, we have that $\push T_{\tilde X}\to T_{X_0}^0$ is an isomorphism and $R^1\push T_{\tilde X}=R^2\push T_{\tilde X}=0.$
  Thus, from the Leray spectral sequence, we have that $H^2(X_0,T_{X_0}^0)\cong H^2(\tilde X, T_{\tilde X})$. 

  Also, we have that $R^1\sigma_* T_{\tilde X}=R^2\sigma_* T_{\tilde X}=0$ and, from the natural map $\sigma_* T_{\tilde X}\to T_X$, there is an exact sequence:
  $$0\to \sigma_* T_{\tilde X}\to T_X\to \mathcal{S}\to 0, $$
  where $S$ is a skyscraper sheaf supported at the blowup points $\xi$. In particular, via Leray and the vanishing of $H^i(X,\mathcal{S})$ for $i>0$, 
  $$H^2(X_0, T_{X_0}^0)\cong H^2(\tilde X, T_{\tilde X})\cong H^2(X, \sigma_* T_{\tilde X})\cong H^2(X, T_X)=0, $$
  since $X$ is Fano and via Nakano vanishing. So finally $\mathbb{T}_{X_0}^2=0$ and the map $\mathbb{T}_{X_0}^1\to H^0(X_0, T_{X_0}^1)$ is surjective. This surjectivity guarantees that any local first-order smoothing lifts to a global first-order deformation. Combined with the vanishing of the obstruction space $\mathbb{T}_{X_0}^2(-\log Y_0) = 0$, these deformations are unconditionally unobstructed and extend to higher orders.
\end{proof}

\subsection{Deformations of the pair: del Pezzo boundary}
\begin{prop}
  Assume that $Y_0$ is a del Pezzo surface. Hence
  $$\Ext^2(\Omega_{X_0}^1(\log Y_0),\OO_{X_0})\cong \mathbb{T}_{X_0}^2(-\log Y_0)=0. $$
\end{prop}
\begin{proof}
  It follows from 
  $$H^1(Y_0, N_{Y_0/X_0})=H^2(Y_0, N_{Y_0/X_0})=0 $$
  by using Kodaira vanishing.

  We conclude the proof by using the exact sequences (\ref{eq:LEStanToNor}) in subsection \ref{ss:setting}, which imply that $H^2(X_0, T_{X_0}^0(-\log Y_0))=0$, from which the result follows easily.
\end{proof}

Using \ref{eq:locToGlo} and \ref{eq:DefPairToDef} we obtain the following Corollary.

\begin{cor}
  The natural map $\mathbb{T}_{X_0}^1(-\log Y_0)\to H^0(T_{X_0}^1)$ is surjective, and the deformations of the pair $(X_0,Y_0)$ are unobstructed.
  
  Moreover, $\mathbb{T}_{X_0}^1(-\log Y_0)\to \mathbb{T}_{X_0}^1$ is surjective and the morphism from the deformation functor of the pair $(X_0,Y_0)$ to the deformation functor of $X_0$ is smooth.
\end{cor}

\begin{rem}
  In general, we don't expect $Y_0$ to be a stable submanifold of $X_0$. For example, a deformation of $\tilde X$ and hence $X_0$ allows for moving the points of $\xi$ arbitrarily. 
  But once $e$ is large enough, there won't be a $Y\in \abs{-\frac{1}{2} K_X}$ passing through $\xi$. This is reflected in the fact that, by Riemann-Roch, $\xi(Y_0, K_{Y_0}^{-1})=c_1(K_{Y_0})^2+1$, so that $h^1(Y_0, K_{Y_0}^{-1})\not = 0$ as soon as $c_1(K_{Y_0})^2\leq -2$.
\end{rem}

\begin{rem}
If we blow-up less points than the degree of $X$, this is $e<d$, we can show that $H^0(-K_{\tilde {X}})$ is non-empty, basepoint free, its generic element $S$ is smooth, disjoint from $\Gamma$ and by Kawamata-Viehweg $H^1(S,N_{S/\tilde X})=0$. Hence the log-Calabi-Yau pair $(X_0,S)$ is also unobstructed.
\end{rem}


\subsection{Deformations of the pair: general case}

\begin{thm}\label{thm:unobstructedness}
  Let $(X_0,Y_0)$ be as in subsection \ref{ss:setting}. Then
  \begin{enumerate}
    \item $\mathbb{T}_{X_0}^2(-\log Y_0)=0$,
    \item the map $\mathbb{T}_{X_0}^1(-\log Y_0)\to H^0(T_{X_0}^1)$ is surjective,
    \item $H^2(\tilde{X}, T_{\tilde X}(-\log \tilde Y))=0.$
  \end{enumerate}
   
\end{thm}

\begin{proof}
By using the Leray spectral sequence as in the proof of Proposition \ref{prop:defAbs} and by diagram (\ref{eq:friedman_diagram}), it suffices to show that $H^2(\tilde{X}, T_{\tilde X}(-\log \tilde Y))=0.$

By considering part of the long exact sequence associated to
$$0\to T_{\tilde X}(-\log \tilde Y)\to T_{\tilde X}\to N_{\tilde Y/ \tilde X}\to 0,$$
and using Proposition \ref{prop:defAbs}, the vanishing $H^2(\tilde{X}, T_{\tilde X}(-\log \tilde Y))=0$ is equivalent to show that 
$$H^1(\tilde X, T_{\tilde X})\to H^1(\tilde Y, N_{\tilde Y/ \tilde X}) $$
is surjective, because $$H^2(\tilde X, T_{\tilde X})\cong H^2(X, \sigma_* T_{\tilde X})\cong H^2(T_X)=0$$
since $X$ is Fano and via Nakano vanishing. 

Denote the ideal sheaf on $X$ (respectively on Y) of the points $\xi$ by $I_{\xi/X}$ (resp. by $I_{\xi/Y}$).
We have isomorphisms
$$H^1(\tilde X, T_{\tilde X})\cong H^1(X, \sigma_* T_{\tilde X})\cong H^1(X, T_X\otimes I_{\xi/X}). $$
using this and the coboundary map $\partial$ given by the following sequence:
$$0\to T_X\otimes I_{\xi/X}\to T_X\to N_{\xi/X}\to 0, $$
we obtain $\partial: H^0(\xi, N_{\xi/X})\to H^1(\tilde X, T_{\tilde X})$. Its image consists of the tangents to deformations on $\tilde X$ arising from keeping $X$ fixed but deforming the blowup points $\xi$.

Let $\bar{\sigma}:=\sigma\vert_{\tilde Y}$ and let $\bar{E}_i =E_i\cap \tilde{Y}$ be the restrictions of the exceptional divisors $E_i$.
Since $N_{\tilde Y/\tilde X}\cong \bar{\sigma}^* N_{Y/X}\otimes \sum_i \OO_{\tilde Y}(-\bar E_i)$, we have $\bar{\sigma}_* N_{\tilde{Y}/\tilde{X}}=N_{Y/X}\otimes I_{\xi/Y}$, and $R^i \bar{\sigma}_* N_{\tilde Y/\tilde X}=0$ for $i>0$.
Twisting the ideal sheaf sequence defining $\xi$ inside $Y$, we obtain a coboundary map:
$$\bar{\partial}:H^0(\xi, N_{Y/X}\otimes \OO_{\xi})\to H^1(Y, N_{Y/X}\otimes I_{\xi/Y})=H^1(Y, \bar{\sigma}_* N_{\tilde Y/\tilde X})=H^1(\tilde Y, N_{\tilde Y/ \tilde X}). $$
In the present case, $H^1(Y, N_{Y/X})=0$ since $Y$ is a del Pezzo surface and $N_{Y/X}=K_{Y}^{-1}$.
Now, we have a commutative diagram 
$$\begin{tikzcd}
  H^0(\xi,N_{\xi/X}) \ar[r] \ar[d, "\partial"] & H^0(\xi, N_{Y/X}\otimes \OO_\xi) \ar[d, "\bar{\partial} "] \\
  H^1(\tilde X, T_{\tilde X}) \ar[r] & H^1(\tilde Y, N_{\tilde Y/\tilde X})
\end{tikzcd}. $$
The top horizontal map comes from the sequence of normal bundles 
$$0\to N_{\xi/Y}\to N_{\xi/X}\vert_{\xi} \to N_{Y/X}\otimes \OO_{\xi}\to 0.$$ 
This sequence consists of skyscraper sheaves supported on $\xi$. Then $H^0(\xi, N_{\xi/X})\to H^0(\xi, N_{Y/X}\otimes \OO_{\xi})$ is surjective and $\bar{\partial}$ is surjective, so $H^1(\tilde X, T_{\tilde X})\to H^1(\tilde Y, N_{\tilde Y/\tilde X})$ is surjective.
\end{proof}


\begin{thebibliography}{99}
\bibitem[]{F86} Friedman, Robert. Simultaneous resolution of threefold double points. Math. Ann. 274 (1986) 671--689
\end{thebibliography}
\end{document}
